\scanpage{404}\mspage{\deleted{\unsure{14}}143}
Dans $1^\circ$ et $3^\circ$, on veut prouver que certaines flèches
\inserted{($p$ et $\bar f$)} sont à la fois \uline{propres} et \uline{lisses} -- ce qui est vraiment
\unsure{idéal|idiot} ! Je dis qu'il suffit à chaque fois de
prouver l'un des deux, que l'autre s'en déduit
(comme ``cas particulier'') par dualité. Pour ceci,
notons
\[
\tau\unsure{\mapstochar\relbar|\mapsto}\tau^{-}
\]
l'automorphisme (\uline{pas} \uline{anti}-automorphisme) de
$\unsure{\underline{\mathrm{Ty}}|\underline{\mathrm{Tz}}}$ qui consiste à changer de \unsure{signe|signes} toutes les
flèches de transition -- c'est donc, \inserted{au niveau des objets,} l'unique
autom. du monoïde libre $\mathrm{Ty}$, qui échange
$\Delta_+$ et $\Delta_-$. Au système projectif $\underline{\unsure{t|f}}$ de $\underline{\mathrm{Ty}}$
correspond le système projectif $\underline t^{-}$. Ceci
dit, on a un isom. de catégories
\[
\boxed{\underline{\mathrm{Ch}}_{\infty}^{\underline t}(X^{\circ})
\simeq\bigl(\underline{\mathrm{Ch}}_{\infty}^{\underline t^{-}}(X)\bigr)^{\circ}}
\]
% slanted
\margin{cf page 140 pour le passage de $\underline t$ à $\underline t^{\circ}$}
et \deleted{le foncteur} le diagramme commutatif
\diagram{\begin{tikzcd}[ampersand replacement=\&,row sep=large]
\underline{\mathrm{Ch}}_{\infty}^{\underline t}(X^{\circ}) \arrow[d,"p^{\underline t}_{X^{\circ}}"] \& \simeq \& \underline{\mathrm{Ch}}_{\infty}^{\underline t^{-}}(X)^{\circ}\arrow[d,"{(p_X^{\underline t^{-}})^{\circ}}"]\\
\deleted{\unsure{C}}X^{\circ}\times X^{\circ} \& = \& (X\times X)^{\circ}
\end{tikzcd}}
Ainsi, dire que $p^{\underline t}_X$ est lisse, i.e que
$(p^{\underline t}_X)^{\circ}$ est propre, équivaut à dire que $p^{\underline t}_{X^{\circ}}$ est
propre.
% archive sheet number: 403
